\documentclass[11pt, a4paper]{article}

% Gemeinsame Style-Definitionen (TEXINPUTS muss templates/ enthalten — siehe scripts/build_tex.py)
% bfw-style.tex — gemeinsame Style-Definitionen für alle Handouts/Aufgabenhefte
% Voraussetzung: Kompilation mit xelatex (wegen fontspec)

\usepackage[ngerman]{babel}
\usepackage{geometry}
\geometry{a4paper, top=2.2cm, bottom=2.2cm, left=2.3cm, right=2.3cm}

\usepackage{fontspec}
% Fallback-Kette: Source Sans Pro -> Calibri -> Segoe UI -> Latin Modern Sans
% (Latin Modern ist immer mit MiKTeX vorhanden)
\IfFontExistsTF{Source Sans Pro}{%
  \setmainfont{Source Sans Pro}[Ligatures=TeX]%
}{\IfFontExistsTF{Calibri}{%
  \setmainfont{Calibri}[Ligatures=TeX]%
}{\IfFontExistsTF{Segoe UI}{%
  \setmainfont{Segoe UI}[Ligatures=TeX]%
}{%
  \setmainfont{Latin Modern Sans}[Ligatures=TeX]%
}}}
\IfFontExistsTF{Source Code Pro}{%
  \setmonofont{Source Code Pro}[Scale=0.92]%
}{\IfFontExistsTF{Consolas}{%
  \setmonofont{Consolas}[Scale=0.92]%
}{%
  \setmonofont{Latin Modern Mono}[Scale=0.92]%
}}

\usepackage{xcolor}
% Farbpalette: hell, freundlich, modern
\definecolor{bfwPrimary}{HTML}{0EA5A4}    % Petrol/Türkis - Primär
\definecolor{bfwPrimaryDark}{HTML}{0F766E} % dunkler Petrol für Headings
\definecolor{bfwAccent}{HTML}{F59E0B}     % Amber - Akzent
\definecolor{bfwSuccess}{HTML}{10B981}    % Grün - Erfolg / Lösung
\definecolor{bfwWarn}{HTML}{F97316}       % Orange - Achtung
\definecolor{bfwInk}{HTML}{0F172A}        % fast Schwarz
\definecolor{bfwInkSoft}{HTML}{475569}    % gedämpftes Grau
\definecolor{bfwBgSoft}{HTML}{F8FAFC}     % off-white
\definecolor{bfwBgTeal}{HTML}{ECFEFF}     % helles Türkis
\definecolor{bfwBgAmber}{HTML}{FEF3C7}    % helles Gelb
\definecolor{bfwBgRose}{HTML}{FFE4E6}     % helles Rosé für Eselsbrücken

\usepackage[colorlinks=true, linkcolor=bfwPrimaryDark, urlcolor=bfwPrimaryDark]{hyperref}
\usepackage{enumitem}
\setlist[itemize]{leftmargin=*, itemsep=2pt, topsep=2pt}
\setlist[enumerate]{leftmargin=*, itemsep=2pt, topsep=2pt}

\usepackage[most]{tcolorbox}
\usepackage{booktabs}
\usepackage{tikz}
\usetikzlibrary{decorations.pathmorphing, positioning, shapes.geometric, arrows.meta}

% Merkkästchen — wichtige Definition / Kernpunkt
\newtcolorbox{merksatz}[1][]{%
  enhanced, breakable,
  colback=bfwBgTeal,
  colframe=bfwPrimary,
  boxrule=0.6pt,
  arc=4pt,
  left=10pt, right=10pt, top=8pt, bottom=8pt,
  fonttitle=\bfseries\color{bfwPrimaryDark},
  title={\faLightbulb\ Merksatz},
  #1
}

% Eselsbrücke — Mnemoniks
\newtcolorbox{eselsbruecke}[1][]{%
  enhanced, breakable,
  colback=bfwBgRose,
  colframe=bfwAccent,
  boxrule=0.6pt,
  arc=4pt,
  left=10pt, right=10pt, top=8pt, bottom=8pt,
  fonttitle=\bfseries\color{bfwWarn},
  title={Eselsbrücke},
  #1
}

% Beispiel-Box
\newtcolorbox{beispiel}[1][]{%
  enhanced, breakable,
  colback=bfwBgSoft,
  colframe=bfwInkSoft,
  boxrule=0.4pt,
  arc=3pt,
  left=10pt, right=10pt, top=8pt, bottom=8pt,
  fonttitle=\bfseries\color{bfwInk},
  title={Beispiel},
  #1
}

% Lösungs-Box (für Aufgabenhefte)
\newtcolorbox{loesung}[1][]{%
  enhanced, breakable,
  colback=bfwBgAmber!40,
  colframe=bfwSuccess,
  boxrule=0.6pt,
  arc=3pt,
  left=10pt, right=10pt, top=8pt, bottom=8pt,
  fonttitle=\bfseries\color{bfwSuccess!70!black},
  title={Lösung},
  #1
}

% Glossar-Eintrag
\newcommand{\glossareintrag}[2]{%
  \par\noindent\textbf{\color{bfwPrimaryDark}#1} \enspace #2\par\smallskip
}

% Section-Stil — etwas farbiger als Standard
\usepackage{titlesec}
\titleformat{\section}
  {\Large\bfseries\color{bfwPrimaryDark}}
  {\thesection}{0.6em}{}
\titleformat{\subsection}
  {\large\bfseries\color{bfwPrimary}}
  {\thesubsection}{0.5em}{}
\titleformat{\subsubsection}
  {\normalsize\bfseries\color{bfwInk}}
  {\thesubsubsection}{0.4em}{}

% TikZ-Stil "skizzenhaft" — etwas weicher als Rough.js, aber stabil
\tikzset{
  roughbox/.style = {
    draw=bfwPrimaryDark, thick, fill=bfwBgTeal,
    rounded corners=4pt, inner sep=6pt
  },
  roughaccent/.style = {
    draw=bfwAccent, thick, fill=bfwBgAmber,
    rounded corners=4pt, inner sep=6pt
  },
  roughline/.style = {
    draw=bfwInkSoft, thick, -{Stealth[length=2.5mm]}
  }
}

% Bullet-Symbole (bewusst kein FontAwesome — vermeidet Font-Installations-Probleme)
\newcommand{\faLightbulb}{$\bullet$}
\newcommand{\faBookmark}{$\bullet$}

% Header/Footer
\usepackage{fancyhdr}
\pagestyle{fancy}
\fancyhf{}
\renewcommand{\headrulewidth}{0pt}
\fancyfoot[L]{\small\color{bfwInkSoft}\thepage}
\fancyfoot[R]{\small\color{bfwInkSoft} BFW M-064 Beschaffung im Büromanagement}


\title{\vspace*{-1.5cm}\Huge\bfseries\color{bfwPrimaryDark}Tag~14: Digitales DMS, Datenschutz \& Umweltschutz\\[0.4em]
  \large\color{bfwInkSoft}\textnormal{Dokumente digital verwalten, Daten schützen, nachhaltig wirtschaften}}
\author{\textsf{Büroprozesse gestalten und Arbeitsvorgänge organisieren \textperiodcentered{} M-067}\\
\textsf{\small\copyright\ Can Siebert\,\textperiodcentered\,2026}}
\date{16.07.2026}

\begin{document}
\maketitle
\thispagestyle{empty}

\vspace{1em}
\begin{merksatz}[title={\bfseries Worum geht es heute?}]
Modernes Büromanagement bedeutet, mit Informationen verantwortungsvoll umzugehen.
Drei große Themen greifen dabei ineinander: Wir verwalten Dokumente \emph{digital}
in einem Dokumentenmanagementsystem (DMS), wir \emph{schützen} die darin enthaltenen
Daten nach den Vorgaben der EU-DSGVO und des BDSG, und wir handeln \emph{nachhaltig},
indem wir Umweltgesetze beachten und Ressourcen schonen. In diesem Handout lernen Sie,
wie revisionssichere Archivierung funktioniert, wie sich Datenschutz und Datensicherheit
unterscheiden und welche gesetzlichen Grundlagen den betrieblichen Umweltschutz prägen.
Das ist sowohl Berufsalltag als auch fester Bestandteil der IHK-Abschlussprüfung.
\end{merksatz}

\tableofcontents
\vspace{1em}
\hrule
\vspace{1em}

\section{Elektronische Dokumente und das DMS}

Durch die zunehmende Nutzung elektronischer Medien stellt sich für jedes Unternehmen
die Frage, wie elektronische Dokumente gespeichert und verwaltet werden. Wird kein
Datenmanagementsystem genutzt, sollten Texte und Daten zumindest nach einem
\emph{einheitlichen Dateiablagesystem} gespeichert werden: Die Ordnerstrukturen der
elektronischen Ablage und der Papierablage sollten aufeinander abgestimmt und die
Dateibenennung einheitlich geregelt sein.

Als \emph{Datenträger} kommen für die magnetische Speicherung Festplatten und
Magnetbänder zum Einsatz, als optische Speicher CD und DVD; mobil wird auf Speicherkarten
und USB-Sticks gespeichert. Für große Datenmengen mit sehr langer Lebensdauer (zwischen
30 und 100 Jahren) eignen sich auch heute noch Mikrofilme.

Unternehmen mit hohem Schriftgutaufkommen nutzen zunehmend \emph{elektronische
Dokumentenmanagementsysteme (DMS)}. Dabei werden ursprünglich papiergebundene
Schriftstücke in elektronische Dokumente umgewandelt und datenbankgestützt digital
archiviert.

\begin{merksatz}
Ein DMS wandelt Papierdokumente in elektronische Dokumente um und verwaltet sie
zentral und datenbankgestützt. Es ist die technische Grundlage des \emph{papierlosen
Büros}.
\end{merksatz}

\subsection{Vorteile eines DMS}

Ein DMS bietet einem Unternehmen unter anderem folgende Vorteile:

\begin{itemize}
  \item schneller Zugriff auf Dokumente und Informationen (Volltextsuche),
  \item Wegfall von Verteilerkopien,
  \item der Verlust von Dokumenten wird reduziert,
  \item geringer Platzbedarf für die Archivierung,
  \item weitgehender Wegfall monotoner Ablagetätigkeiten,
  \item unabhängiger Zugriff auf Dokumente, auch wenn diese gerade in Bearbeitung sind.
\end{itemize}

\subsection{Revisionssichere Archivierung, Workflow und Cloud}

Der Gesetzgeber gibt für Unterlagen, für die eine \emph{Aufbewahrungsfrist} vorgesehen
ist, die Regelung für eine elektronische Datenspeicherung vor (Auszug aus dem
HGB \S~257~(3)). Damit ein DMS aufbewahrungspflichtige Unterlagen ersetzen darf, muss
es \emph{revisionssicher} arbeiten. Die Anforderungen konkretisieren die \emph{GoBD}
(Grundsätze zur ordnungsmäßigen Führung und Aufbewahrung von Büchern, Aufzeichnungen
und Unterlagen in elektronischer Form).

\begin{merksatz}[title={\bfseries Revisionssicher heißt}]
Gespeicherte Dokumente sind \emph{unveränderbar}, \emph{vollständig}, \emph{nachvollziehbar}
(Protokollierung) und über die gesamte Aufbewahrungsfrist \emph{wiederauffindbar}.
\end{merksatz}

Moderne DMS bieten darüber hinaus:
\begin{itemize}
  \item \textbf{Workflow:} Dokumente durchlaufen automatisch definierte Bearbeitungs- und Freigabeschritte (z.\,B.\ Rechnungsprüfung, Freigabe, Buchung).
  \item \textbf{Versionierung:} jede Änderung wird als neue Version gespeichert; ältere Stände bleiben nachvollziehbar erhalten.
  \item \textbf{Cloud-DMS:} Das System wird als Online-Dienst genutzt; der Zugriff ist ortsunabhängig (wichtig für Homeoffice und mobiles Arbeiten), erfordert aber besondere Aufmerksamkeit für Datenschutz und Vertragsgestaltung.
\end{itemize}

\begin{eselsbruecke}
\textbf{„Unveränderbar, Vollständig, Nachvollziehbar, Wiederauffindbar”} --- mit
\emph{U-V-N-W} merken Sie sich die vier Kernanforderungen der revisionssicheren
Archivierung.
\end{eselsbruecke}

\section{Datenschutz und Datensicherheit}

Jedes Unternehmen ist verpflichtet, für die Sicherheit der gespeicherten Daten zu
sorgen. Durch geeignete Maßnahmen lassen sich Bedienungsfehler, technische Störungen,
Datenverluste und Datenmissbrauch verhindern oder ausschließen. Die beiden zentralen
Begriffe sind dabei sauber zu unterscheiden:

\begin{center}
\begin{tabular}{p{2.6cm} p{5.1cm} p{4.4cm}}
\toprule
 & \textbf{Datenschutz} & \textbf{Datensicherheit}\\
\midrule
Geschützt wird & natürliche Personen & Hard-, Software, Daten\\
Gefahr & Verletzung von Persönlichkeitsrechten & Verlust, Zerstörung, Missbrauch durch Unbefugte\\
\bottomrule
\end{tabular}
\end{center}

\medskip
Alle Maßnahmen erfolgen unter Beachtung der \emph{EU-Datenschutz-Grundverordnung
(EU-DSGVO)} und des \emph{Bundesdatenschutzgesetzes (BDSG)}.

\begin{eselsbruecke}
\textbf{Daten\underline{schutz} schützt Menschen, Daten\underline{sicherheit} sichert
Technik.} Wer schützt wen --- das ist die häufigste Verwechslung in der Prüfung.
\end{eselsbruecke}

\subsection{Personenbezogene Daten und Rechtsgrundlagen}

Der \textbf{Datenschutz} beschreibt die Maßnahmen, die ergriffen werden, um personen-
und unternehmensbezogene Daten vor Diebstahl und Missbrauch zu schützen.
\emph{Personenbezogene Daten} sind nach Art.~4 Nr.~1 EU-DSGVO „alle Informationen, die
sich auf eine identifizierte oder identifizierbare natürliche Person beziehen”. Eine
Person ist z.\,B.\ identifizierbar über den Namen, die E-Mail-Adresse, Bankdaten, die
Kreditkartennummer oder die IP-Adresse.

\subsection{Der Datenschutzbeauftragte}

Unternehmen müssen einen \emph{Datenschutzbeauftragten} ernennen, wenn ihr Kerngeschäft
die Überwachung personenbezogener Daten ist (Art.~37 EU-DSGVO). Ergänzend gilt nach
\S~38 BDSG (in der seit 2019 geltenden Fassung): Ein Datenschutzbeauftragter ist zu
ernennen, wenn mindestens \emph{zwanzig Personen} ständig mit der automatisierten
Verarbeitung personenbezogener Daten beschäftigt sind.

Der Datenschutzbeauftragte hat eine Kontroll- und Beratungsfunktion. Er ist der
Unternehmensleitung unmittelbar unterstellt, hinsichtlich seiner Kontrollaufgaben aber
weisungsfrei und genießt besonderen Kündigungsschutz. Wichtig: Der Leiter der
IT-Abteilung kann in der Regel \emph{nicht} Datenschutzbeauftragter werden, da er sich
sonst in bestimmten Bereichen selbst kontrollieren müsste (Interessenkonflikt). Zu
seinen Aufgaben gehören:

\begin{itemize}
  \item Unterrichtung und Beratung der Unternehmensleitung zu allen Datenschutzvorschriften,
  \item Überwachung der Einhaltung der Datenschutzvorschriften,
  \item Sensibilisierung und Schulung der mit der Verarbeitung beschäftigten Mitarbeiter,
  \item Zusammenarbeit mit der Aufsichtsbehörde,
  \item Beratung bei der Datenschutz-Folgenabschätzung.
\end{itemize}

\subsection{Rechte der Betroffenen}

\emph{Betroffene} sind die Personen, deren Daten erhoben, verarbeitet und genutzt werden.
Die EU-DSGVO gewährt ihnen folgende Rechte:

\begin{center}
\begin{tabular}{p{5.2cm} p{6.6cm}}
\toprule
\textbf{Recht (Artikel)} & \textbf{Bedeutung}\\
\midrule
Auskunftsrecht (Art.~15) & Information über Art, Empfänger und Zweck der gespeicherten Daten\\
Recht auf Berichtigung (Art.~16) & Korrektur falscher personenbezogener Daten\\
Recht auf Löschung (Art.~17) & „Recht auf Vergessenwerden” bei unzulässiger Speicherung\\
Recht auf Einschränkung (Art.~18) & Verarbeitung wird während einer Prüfung ausgesetzt\\
Recht auf Datenübertragbarkeit (Art.~20) & Herausgabe der Daten in maschinenlesbarem Format\\
Widerspruchsrecht (Art.~21) & Widerspruch gegen die Verarbeitung jederzeit möglich\\
\bottomrule
\end{tabular}
\end{center}

\subsection{Weitere Maßnahmen und Aufbewahrung vs.\ Löschpflicht}

Neben dem Datenschutzbeauftragten sind weitere \emph{technisch-organisatorische
Maßnahmen} nötig: passwortgeschützte Benutzerkonten, Antivirenschutz auf jedem Rechner
im Netzwerk, Firewalls zum Schutz vor unerlaubten Zugriffen sowie Sensibilisierung und
Schulung der Mitarbeiter (häufig durch eine zu unterzeichnende Datenschutzerklärung).

Zum Datenschutz gehört auch der Umgang mit Unterlagen, deren Aufbewahrungsfrist
abgelaufen ist: Diese Schriftstücke \emph{dürfen} (und müssen aus Datenschutzgründen)
vernichtet werden. Eine sichere Vernichtung gelingt mit \emph{Aktenvernichtern}
(Streifenschnitt oder --- sicherer --- Partikelschnitt).

\begin{beispiel}
\textbf{Aufbewahrung vs.\ Löschpflicht:} Eine Rechnung muss handels- und steuerrechtlich
mehrere Jahre aufbewahrt werden --- so lange darf sie \emph{nicht} gelöscht werden.
Ist die Frist abgelaufen und werden die Daten nicht mehr benötigt, greift dagegen die
\emph{Löschpflicht} aus dem Datenschutz: Die Unterlagen sind dann zu vernichten.
\end{beispiel}

\section{Datensicherung}

Die \textbf{Datensicherung} umfasst alle Maßnahmen zum Schutz von Daten, Datenträgern
und Programmen vor Übertragungsfehlern, Verfälschungen, Vernichtung, Diebstahl und
unberechtigtem Zugriff. Sie ist in fünf Bereichen zu gewährleisten:

\begin{center}
\begin{tabular}{p{3.4cm} p{8.4cm}}
\toprule
\textbf{Gefahrenbereich} & \textbf{Maßnahmen}\\
\midrule
Datenerfassung & Prüfziffernverfahren, Plausibilitätskontrolle\\
Datenübertragung & Prüfbit (Paritätsbit), Übertragungsprotokolle, Verschlüsselung (Kryptographie)\\
Datenträger & Schreibschutz, Backup nach dem Generationenprinzip, Archivierung (feuerfester Schrank)\\
Rechnerausfall & Notstromanlage (USV), Plattenspiegelung (Duplexing)\\
Computerkriminalität & Zugriffskontrolle, Passwort, Kennzahl (PIN), Virenschutzprogramme\\
\bottomrule
\end{tabular}
\end{center}

\subsection{Backup nach dem Generationenprinzip}

Von jedem Datenbestand müssen immer mindestens \emph{drei Generationen} aufeinander
aufbauender Datensicherungen (Backups) vorhanden sein. Die älteste Sicherung heißt
\emph{Großvater}, die zweitälteste \emph{Vater} und die jüngste \emph{Sohn}.

\begin{beispiel}
Montags werden alle Daten auf Band~1 gesichert (\glqq Sohn\grqq). Dienstags werden die
Daten auf Band~2 gesichert (\glqq Sohn\grqq); Band~1 wird zum \glqq Vater\grqq.
Mittwochs werden alle Daten auf Band~3 gesichert (\glqq Sohn\grqq); Band~2 wird
\glqq Vater\grqq{} und Band~1 wird \glqq Großvater\grqq.
\end{beispiel}

Eine \emph{USV} (unterbrechungsfreie Stromversorgung) wird zwischen Computer und
Steckdose geschaltet und versorgt den Rechner bei Stromausfall so lange weiter, bis
die Daten aus dem Hauptspeicher gesichert sind.

\section{Umweltschutz im Unternehmen}

Alle Maßnahmen zum Schutz der Umwelt und zur Vermeidung von Umweltbelastungen werden
als \emph{Umweltschutz} bezeichnet. Ein aktiver Umweltschutz wirkt sich oft auch
wirtschaftlich positiv aus --- etwa wenn ein Unternehmen durch den Einsatz von
Recyclingpapier zugleich die Umwelt entlastet und Kosten spart.

\subsection{Gesetzliche Grundlagen}

\subsubsection{Kreislaufwirtschaftsgesetz (KrWG)}

Das \textbf{Kreislaufwirtschaftsgesetz (KrWG)} ist das Bundesgesetz zum Abfallrecht.
Ziel ist es, Abfälle zu reduzieren und natürliche Ressourcen zu schonen. \S~6 KrWG legt
die \emph{Abfallhierarchie} fest --- Vorrang hat stets die aus Umweltsicht beste
Möglichkeit (\emph{Vermeidung vor Verwertung, Verwertung vor Beseitigung}):

\begin{enumerate}
  \item \textbf{Vermeidung} --- Abfälle in Menge und Schädlichkeit reduzieren (z.\,B.\ Mehrwegverpackungen, Nachfüllpackungen).
  \item \textbf{Vorbereitung zur Wiederverwendung} --- Prüfung, Reinigung, Reparatur, damit Erzeugnisse erneut zum gleichen Zweck genutzt werden.
  \item \textbf{Recycling} --- gebrauchte Materialien werden zum Rohstoff (Sekundärrohstoff) wiederverwertet (z.\,B.\ Altpapier zu Recyclingpapier).
  \item \textbf{Sonstige Verwertung} --- insbesondere energetische Verwertung (Verbrennung mit Energiegewinnung).
  \item \textbf{Beseitigung} --- z.\,B.\ Ablagerung auf Deponien.
\end{enumerate}

Zu unterscheiden sind: \emph{Wiederverwertung} (Altstoffe kommen im gleichartigen Prozess
erneut zum Einsatz, z.\,B.\ Glas einschmelzen), \emph{Wiederverwendung} (ein gebrauchtes
Produkt wird für denselben Zweck genutzt, z.\,B.\ Glaspfandflaschen) und
\emph{Weiterverwendung} (ein gebrauchtes Produkt wird für einen anderen Zweck genutzt,
z.\,B.\ alte Autoreifen als Schiffsfender).

\subsubsection{Verpackungsgesetz (VerpackG)}

Das \textbf{Verpackungsgesetz (VerpackG)} trat zum 1.~Januar 2019 in Kraft und löste die
Verpackungsverordnung ab. Hersteller müssen sich bei der \emph{Zentralen Stelle
Verpackungsregister} registrieren und sich an einem \emph{dualen Entsorgungssystem}
(z.\,B.\ Grüner Punkt --- DSD) beteiligen, um für die Entsorgungskosten aufzukommen.
Das duale System umfasst die drei Säulen \emph{Altpapier, Glas} sowie \emph{gelber
Sack/gelbe Tonne}. Das VerpackG unterscheidet in \S~3 vier Verpackungstypen:

\begin{itemize}
  \item \textbf{Verkaufsverpackungen} (Bestandteil der Verkaufseinheit; auch Service- und Versandverpackungen),
  \item \textbf{Umverpackungen} (enthalten mehrere Verkaufseinheiten),
  \item \textbf{Transportverpackungen} (erleichtern Handhabung und Transport, z.\,B.\ Gitterboxen, Euro-Paletten).
\end{itemize}

\subsubsection{Elektro- und Elektronikgerätegesetz (ElektroG)}

Das \textbf{Elektro- und Elektronikgerätegesetz (ElektroG)} regelt die umweltverträgliche
Entsorgung von Elektro- und Elektronikgeräten. Private Haushalte dürfen Altgeräte
\emph{nicht} über den Restmüll entsorgen; sie sind durch das Symbol der
\emph{durchgestrichenen Mülltonne} gekennzeichnet. Hersteller müssen sich vor dem
Inverkehrbringen bei der \glqq Gemeinsamen Stelle der Hersteller\grqq{} (Stiftung EAR)
registrieren --- bei Verstößen drohen Bußgelder bis zu 100\,000,00~Euro. Große Händler
(Verkaufs- bzw.\ Lagerfläche über 400~m\textsuperscript{2}) sind zur Rücknahme
verpflichtet; kleine Geräte (Kantenlänge unter 25~cm) müssen sie immer kostenlos
zurücknehmen.

\subsection{Nachhaltiges Wirtschaften}

Der Begriff \emph{Nachhaltigkeit} stammt aus der Forstwirtschaft: Es dürfen nicht mehr
Bäume gefällt werden, als nachwachsen. Übertragen heißt das, heutige Bedürfnisse so zu
befriedigen, dass den nachfolgenden Generationen alle Ressourcen in gleicher Menge und
Qualität zur Verfügung stehen. Zum nachhaltigen Wirtschaften zählen umweltfreundliche
Produkte, die Reduzierung des Energieverbrauchs, umweltschonende Produktionsverfahren,
rationelle Ressourcenverwendung und fairer Handel (\emph{Fair Trade}).

\subsubsection{Umweltfreundliche Produkte und Siegel}

Beim Einkauf helfen \emph{Umweltprüfsiegel}, umweltfreundliche Produkte zu erkennen:

\begin{center}
\begin{tabular}{p{3.6cm} p{8.2cm}}
\toprule
\textbf{Siegel} & \textbf{Bedeutung}\\
\midrule
Blauer Engel & Umweltzeichen der Bundesregierung für umwelt- und gesundheitsfreundliche Produkte\\
Energy Star & Kennzeichnung stromsparender IT-Geräte (US-Umweltbehörde EPA)\\
EU-Umweltblume (Ecolabel) & EU-Siegel über den gesamten Lebenszyklus eines Produkts\\
TCO Certified & nachhaltige Gestaltung, Herstellung und Verwendung von IT-Produkten\\
FSC / PEFC & Holz- und Papierprodukte aus nachhaltig bewirtschafteten Wäldern\\
\bottomrule
\end{tabular}
\end{center}

\subsubsection{Energie sparen und Green IT}

Einen großen Beitrag zum Klima- und Umweltschutz leistet die Verminderung des
Energieverbrauchs. Im Büro helfen z.\,B.: Licht und Bildschirm beim Verlassen des Raums
abschalten, den Drucker erst bei Bedarf einschalten, im Winter stoßlüften statt
Kippstellung, Verkehrsflächen mit bewegungsgesteuertem Licht ausstatten, Geräte zum
Feierabend komplett abschalten (Stand-by verbraucht weiter Strom!), auf LED umstellen
und bei Neuanschaffungen auf Energieeffizienz achten.

\emph{Green IT} bezeichnet den umwelt- und ressourcenschonenden Einsatz von
Informationstechnik über den gesamten Lebenszyklus --- etwa durch energieeffiziente
Server, Virtualisierung statt vieler Einzelrechner oder die bewusste Nutzung von
Cloud-Diensten. Zusammen mit dem \emph{papierlosen Büro} (digitale Dokumente im DMS)
schließt sich hier der Kreis zum ersten Themenblock dieses Tages.

\section{Zusammenfassung}

\begin{enumerate}
  \item Ein \textbf{DMS} wandelt Papier in elektronische Dokumente um und verwaltet sie datenbankgestützt --- mit schnellem Zugriff, weniger Verlust und geringem Platzbedarf.
  \item \textbf{Revisionssichere Archivierung} (HGB \S~257, GoBD) verlangt unveränderbare, vollständige, nachvollziehbare und wiederauffindbare Dokumente; Workflow, Versionierung und Cloud-DMS ergänzen moderne Systeme.
  \item \textbf{Datenschutz} schützt natürliche Personen, \textbf{Datensicherheit} schützt Hard-, Software und Daten --- Grundlagen sind EU-DSGVO und BDSG.
  \item Der \textbf{Datenschutzbeauftragte} (ab 20 Beschäftigten mit automatisierter Verarbeitung) kontrolliert und berät; Betroffene haben Rechte nach Art.~15--21 DSGVO.
  \item \textbf{Datensicherung} sichert Daten in fünf Bereichen ab; das \textbf{Generationenprinzip} (Großvater--Vater--Sohn) verlangt mindestens drei Backup-Generationen.
  \item \textbf{Umweltschutz}: KrWG (Abfallhierarchie), VerpackG (duales System) und ElektroG bilden den Rahmen; \textbf{nachhaltiges Wirtschaften} setzt auf Umweltsiegel, Energiesparen, papierloses Büro und Green IT.
\end{enumerate}

\newpage

\section*{Glossar}
\addcontentsline{toc}{section}{Glossar}

\glossareintrag{DMS}{Dokumentenmanagementsystem; wandelt papiergebundene Schriftstücke in elektronische Dokumente um und verwaltet sie datenbankgestützt.}

\glossareintrag{Revisionssichere Archivierung}{Speicherung, bei der Dokumente unveränderbar, vollständig, nachvollziehbar und über die Aufbewahrungsfrist wiederauffindbar sind (HGB \S~257, GoBD).}

\glossareintrag{GoBD}{Grundsätze zur ordnungsmäßigen Führung und Aufbewahrung von Büchern, Aufzeichnungen und Unterlagen in elektronischer Form.}

\glossareintrag{Workflow}{Automatisch gesteuerter Ablauf von Bearbeitungs- und Freigabeschritten eines Dokuments im DMS.}

\glossareintrag{Versionierung}{Speicherung jeder Änderung als eigene Version, sodass ältere Stände nachvollziehbar bleiben.}

\glossareintrag{Cloud-DMS}{Dokumentenmanagement als Online-Dienst mit ortsunabhängigem Zugriff.}

\glossareintrag{Datenschutz}{Maßnahmen zum Schutz natürlicher Personen vor Verletzung ihrer Persönlichkeitsrechte durch Missbrauch personenbezogener Daten.}

\glossareintrag{Datensicherheit}{Maßnahmen zum Schutz von Hard-, Software und Daten vor Verlust, Zerstörung und Missbrauch.}

\glossareintrag{Personenbezogene Daten}{Alle Informationen, die sich auf eine identifizierte oder identifizierbare natürliche Person beziehen (Art.~4 Nr.~1 EU-DSGVO).}

\glossareintrag{EU-DSGVO}{Europäische Datenschutz-Grundverordnung; zentrale Rechtsgrundlage des Datenschutzes.}

\glossareintrag{BDSG}{Bundesdatenschutzgesetz; ergänzt die EU-DSGVO auf nationaler Ebene.}

\glossareintrag{Datenschutzbeauftragter}{Person mit Kontroll- und Beratungsfunktion; verpflichtend ab 20 ständig mit automatisierter Datenverarbeitung Beschäftigten (\S~38 BDSG).}

\glossareintrag{Generationenprinzip}{Backup-Verfahren mit mindestens drei aufeinander aufbauenden Sicherungen (Großvater--Vater--Sohn).}

\glossareintrag{USV}{Unterbrechungsfreie Stromversorgung; überbrückt einen Stromausfall, bis Daten gesichert sind.}

\glossareintrag{Kreislaufwirtschaftsgesetz}{Bundesgesetz zum Abfallrecht; legt in \S~6 die Abfallhierarchie fest.}

\glossareintrag{Abfallhierarchie}{Rangfolge Vermeidung, Vorbereitung zur Wiederverwendung, Recycling, sonstige Verwertung, Beseitigung.}

\glossareintrag{Verpackungsgesetz}{Regelt seit 2019 das Inverkehrbringen, die Rücknahme und Verwertung von Verpackungen (duales System).}

\glossareintrag{ElektroG}{Elektro- und Elektronikgerätegesetz; regelt die umweltverträgliche Entsorgung von Elektrogeräten.}

\glossareintrag{Nachhaltigkeit}{Grundsatz, heutige Bedürfnisse so zu befriedigen, dass künftige Generationen über dieselben Ressourcen verfügen.}

\glossareintrag{Green IT}{Umwelt- und ressourcenschonender Einsatz von Informationstechnik über den gesamten Lebenszyklus.}

\end{document}
